\documentclass[11pt]{article}

\usepackage{../ppackage}
\usepackage{bbm}
\usepackage{import}






\usepackage{tikz}
\usetikzlibrary{arrows.meta,patterns}
\usetikzlibrary{ipe} % ipe compatibility library

\usepackage{../../../Notes/tikzit}
\usepackage{../../../Notes/utf8math}

\input{../../../Notes/TIKZ/digraph.tikzstyles}



\usepackage{url}


\newcommand{\solution}{
\bigskip\noindent
	\textbf{Solution: \\}}
	


\DeclareMathOperator{\size}{size}
\DeclareMathOperator{\conv}{conv}
\newcommand{\SV}{\mathrm{SV}}
\newcommand{\bigO}{O}
\newcommand{\cut}{\mathrm{cut}}
\newcommand{\LLL}{\mathrm{LLL}}
\newcommand{\setR}{\mathbb{R}}
\newcommand{\setZ}{\mathbb{Z}}
\newcommand{\setQ}{\mathbb{Q}}
\newcommand{\setC}{\mathbb{C}}
\newcommand{\setN}{\mathbb{N}}
\newcommand{\wt}[1]{\widetilde{#1}}
\newcommand{\opt}{{\sc 0/1-opt}\xspace}
\newcommand{\aug}{{\sc 0/1-aug}\xspace}
\newcommand{\psep}{{\sc 0/1-psep}\xspace}
\newcommand{\sep}{{\sc 0/1-sep}\xspace}
\newcommand{\fopt}{{\sc 0/1-testopt\xspace} }

\newcommand{\hpp}{\mathrm{HPP}}
\newcommand{\nodes}{\mathcal{V}}
\newcommand{\vol}{\mathrm{vol}}
\newcommand{\diag}{\mathrm{diag}}
\newcommand{\arcs}{\mathcal{A}}
\newcommand{\edges}{\mathcal{E}}
\newcommand{\paths}{\mathscr{P}}
\newcommand{\cycles}{\mathcal{C}}




\newcommand{\K}{{\mathcal K}}
\newcommand{\A}{{A}}
\newcommand{\B}{{B}}
\newcommand{\T}{\mathscr{T}}
\newcommand{\eE}{\mathscr{E}}
\newcommand{\eS}{\mathscr{S}}
\newcommand{\eP}{\mathscr{P}}
\newcommand{\eM}{\mathscr{M}}



\newcommand{\transp}{^{\mathrm{T}}}

\newcommand{\smallmat}[1]{\left( \begin{smallmatrix} #1 \end{smallmatrix}\right)}

\newcommand{\mat}[1]{ \begin{pmatrix} #1 \end{pmatrix}}
\newcommand{\smat}[1]{ \big(\begin{smallmatrix} #1 \end{smallmatrix}\big)}

\newcommand{\pc}{\mathscr{P}}
\newcommand{\ob}{\mathscr{O}}
\newcommand{\odds}{\mathscr{W}}
\newcommand{\up}{\mathscr{U}}
\newcommand{\ef}{\mathscr{F}}
\newcommand{\eh}{\mathscr{H}}
\newcommand{\ev}{\mathscr{V}}
\newcommand{\ec}{\mathscr{C}}
\newcommand{\eu}{\mathscr{U}}

\newcommand{\lex}{\mathrm{lex}}

\renewcommand{\leq}{\leqslant}
\renewcommand{\geq}{\geqslant}









\newcommand{\linhull}{\mathrm{lin.hull}}
\newcommand{\affhull}{\mathrm{affine.hull}}
\newcommand{\charcone}{\mathrm{char.cone}}
\newcommand{\cone}{\mathrm{cone}}
\newcommand{\rank}{\mathrm{rank}}
\newcommand{\wb}[1]{\overline{#1}}



\usepackage{enumerate}
\usepackage{booktabs}

	  
\institute{\'Ecole Polytechnique F\'ed\'erale de Lausanne}
\lecture{Discrete Optimization}
\faculty{Prof. Friedrich Eisenbrand}
\term{Spring 2026}
\publishdate{May 5, 2026}
\duedate{ }
\problemset{Problem Set~10 (Solutions)}

\begin{document}
\makeheader

\begin{enumerate}[1)]

\item Suppose we are given three $n \times n$ matrices $A, B, C \in \mathbb{Z}^{n \times n}$ and we want to test whether $A \cdot B = C$ holds. 
We could multiply $A$ and $B$ and then compare the result with $C$. This would amount to running time (number of arithmetic operations) of $O(n^3)$ with the standard matrix-multiplication algorithm. We now show how to perform an efficient \emph{randomized test}. Suppose that you can draw a vector $v \in \{0, 1\}^n$ i.i.d. at random in time $O(n)$. The idea is then to compute the product $B \cdot v$ and then the product $A \cdot (B \cdot v)$ and afterwards $C \cdot v$, all in time $O(n^2)$. Show the following:
\begin{enumerate}[i)]
	\item If $A \cdot B \neq C$, then $\Pr(A \cdot (B \cdot v) = C \cdot v) \leq 1/2$.

	\item Let $v_1, \ldots, v_k \in \{0, 1\}^n$ be i.i.d. at random and suppose that $A \cdot B \neq C$.
	The probability of the event: $A \cdot (B \cdot v_i) = C \cdot v_i$ for each $i = 1, \ldots, k$ is bounded by $1/2^k$.

	\item Conclude that there is an algorithm that runs in time $O(k \cdot n^2)$ which tests whether $A \cdot B = C$ holds. The probability that the algorithm gives the wrong result is bounded by $1/2^k$.
\end{enumerate}

\begin{solution}
\begin{enumerate}[i)]
	\item Since $A \cdot B \neq C$, there exists an index $i \in \{1, \ldots, n\}$ such that the $i$-th row of $A \cdot B$ is different from the $i$-th row of $C$. Let $w^T$ be the vector obtained by subtracting the $i$-th row of $C$ from the $i$-th row of $A \cdot B$. Then we have
	\begin{align*}
		\Pr(A \cdot (B \cdot v) = C \cdot v) &= \Pr((A \cdot B - C) \cdot v = 0) \\
		&\leq \Pr(w^T v = 0).
	\end{align*}
	Since $w \in \Z^n$ and $w \neq 0$, there exists an index $j \in \{1, \ldots, n\}$ such that $w_j \neq 0$. The event $w^T v = 0$ can be rewritten as $w_j \cdot v_j + \sum_{\ell \neq j} w_\ell \cdot v_\ell = 0$. For any fixed values of $\{v_\ell\}_{\ell \neq j}$, at most one of the two values $v_j = 0$ and $v_j = 1$ can satisfy $w_j \cdot v_j + \sum_{\ell \neq j} w_\ell \cdot v_\ell = 0$. Hence we have $\Pr(w^T v = 0) \leq 1/2$.
	\item For each $i=1, \ldots, k$, let $E_i$ be the event that $A \cdot (B \cdot v_i) = C \cdot v_i$. Since the events $E_1, \ldots, E_k$ are independent, by using the bound from i) we get
	\begin{align*}
		\Pr\left(\bigcap_{i=1}^k E_i\right) = \prod_{i=1}^k \Pr(E_i) \leq \prod_{i=1}^k \frac{1}{2} = \frac{1}{2^k}.
	\end{align*}
	\item
	\textbf{Algorithm}: First draw $k$ vectors $v_1, \ldots, v_k \in \{0, 1\}^n$ i.i.d. at random and compute $A \cdot (B \cdot v_i)$ and $C \cdot v_i$ for each $i=1, \ldots, k$. 
	The algorithm returns "Yes" if $A \cdot (B \cdot v_i) = C \cdot v_i$ for all $i=1, \ldots, k$ and "No" otherwise.
	
	\textbf{Runtime}: Drawing $k$ random vectors takes time $O(k \cdot n)$. Computing $B \cdot v_i$, $A \cdot (B \cdot v_i)$, and $C \cdot v_i$ for all $i=1, \ldots, k$ takes time $O(k \cdot n^2)$. Finally comparing $A \cdot (B \cdot v_i)$ and $C \cdot v_i$ for all $i=1, \ldots, k$ takes time $O(k \cdot n)$. Hence the overall runtime is $O(k \cdot n^2)$.

	\textbf{Error probability}: If $A \cdot B = C$, then the algorithm always returns "Yes". If $A \cdot B \neq C$, then by ii) the algorithm returns "Yes" with probability at most $1/2^k$.
\end{enumerate}
\end{solution}

\item Recall that in the runtime analysis of the faster matrix multiplication, we have the following recursion: 
\[
	T(n) = 8 \cdot T(n/2) + \Theta(n^2).
\]
Show that this recursion has the solution $T(n) = \Theta(n^3)$.

\begin{solution}
Apply this recursion $k$ times we get
\[
T(n) = 8^k \cdot T(n/2^k) + \Theta\left(\sum_{i=0}^{k-1} 8^i \cdot (n/2^i)^2\right).
\]
When $k = \log_2(n)$, we have 
\begin{align*}
T(n) = 8^{\log_2(n)} \cdot T(1) + \Theta\left(\sum_{i=0}^{\log_2(n)-1} 8^i \cdot (n/2^i)^2\right).
\end{align*}
The first term is $8^{\log_2(n)} \cdot T(1) = n^{\log_2(8)} \cdot \Theta(1) = \Theta(n^3)$. The second term is 
\[
\Theta\left(n^2 \cdot \sum_{i=0}^{\log_2(n)-1} 8^i \cdot (1/2^i)^2\right) = \Theta\left(n^2 \cdot \sum_{i=0}^{\log_2(n)-1} 2^i\right) = \Theta\left(n^2 \cdot 2^{\log_2(n)}\right) = \Theta(n^3).	
\]

One may also conclude directly from the Master Theorem of the divide-and-conquer recurrences.
\end{solution}

\item Suppose that $n=2^\ell$ and $a, b \in \mathbb{N}$ are two $n$-bit integers. Consider the numbers $a_h$ and $a_l$ which are represented by the first $n/2$ bits and the last $n/2$ bits of a respectively. Likewise the numbers $b_h$ and $b_l$ are the numbers represented by the first half and the second half of the bit-representation of $b$.
\begin{enumerate}[i)]
	\item Show that $a = a_h \cdot 2^{n/2} + a_l$ and $b = b_h \cdot 2^{n/2} + b_l$.

	\item Show that $a \cdot b = a_h \cdot b_h \cdot 2^n + (a_h \cdot b_l + a_l \cdot b_h) \cdot 2^{n/2} + a_l \cdot b_l$.		

	\item Conclude very carefully that two $n$-bit numbers can be multiplied by resorting to three multiplications of $n/2$-bit numbers and $O(n)$ basic operations. 

	\item Conclude that two $n$-bit numbers can be computed in time $O(n^{\log_2(3)})$.
\end{enumerate}

\begin{solution}
	\begin{enumerate}[i)]
		\item It suffices to prove the formula for $a$. Let $x_0, \ldots, x_{n-1} \in \{0, 1\}$ be the bit sequence of $a$. 
		\begin{align*}
		a &= \sum_{i=0}^{n-1} x_i \cdot 2^i \\
		&= \sum_{i=0}^{n/2-1} x_i \cdot 2^i + \sum_{i=n/2}^{n-1} x_i \cdot 2^i \\
		&= \sum_{i=0}^{n/2-1} x_i \cdot 2^i + \left(\sum_{i=n/2}^{n-1} x_i \cdot 2^{i-n/2}\right)2^{n/2} \\
		&= a_l + a_h \cdot 2^{n/2}.
		\end{align*}
		\item Using the formulas from i), we have
		\begin{align*}
		a \cdot b &= (a_h \cdot 2^{n/2} + a_l) \cdot (b_h \cdot 2^{n/2} + b_l) \\
		&= a_h \cdot b_h \cdot 2^n + (a_h \cdot b_l + a_l \cdot b_h) \cdot 2^{n/2} + a_l \cdot b_l.
		\end{align*}
		\item By ii), one can compute $a \cdot b$ by computing $a_h \cdot b_h$, $a_h \cdot b_l + a_l \cdot b_h$, and $a_l \cdot b_l$. 
		An efficient way to do that is to first compute $a_h \cdot b_h$ and $a_l \cdot b_l$. Then compute $(a_h + a_l) \cdot (b_h + b_l)$ and subtract $a_h \cdot b_h$ and $a_l \cdot b_l$ from it to get $a_h \cdot b_l + a_l \cdot b_h$.

		\item Let $T(n)$ be the time to multiply two $n$-bit numbers. By iii), we have the recursion
		\[T(n) = 3 \cdot T(n/2) + O(n).\]
		Solve this recursion to get $T(n) = O(n^{\log_2(3)})$.
	\end{enumerate}
\end{solution}

\item Matrix multiplication can be reduced to matrix inversion. Let $K$ be a field and $A,B ∈ K^{n × n}$.
  \begin{enumerate}[i)]
  \item 
  Let $D ∈ K^{3n ×3n}$ be the invertible  matrix
  \begin{displaymath}
    D =
    \begin{pmatrix}
      I & A & 0 \\
      0 & I & B \\
      0 & 0 & I
    \end{pmatrix}. 
  \end{displaymath}
  Show that
  \begin{displaymath}
    D^{-1}  =
    \begin{pmatrix}
      I & -A & AB \\
      0 & I & -B \\
      0 & 0 & I
    \end{pmatrix}. 
  \end{displaymath}
\item If one can invert an $n×n$ matrix with $T(n)$ field operations, where $T(n) ≥ n^2$, then one can multiply two $n×n$ matrices with  $O(T(n))$ field operations. 
\end{enumerate}

\begin{solution}
\begin{enumerate}[i)]
\item Verify by direct computations.
\item To multiply $A$ and $B$, we construct the matrix $D$ as in i) and compute $D^{-1}$. The product $A \cdot B$ is then given by the upper right block of $D^{-1}$. The overall number of field operations is $O(T(n))$.
\end{enumerate}
\end{solution}

\end{enumerate}



  

\end{document}

%%% Local Variables:
%%% mode: latex
%%% TeX-master: t
%%% End:
